% Part: second-order-logic
% Chapter: syntax-and-semantics
% Section: introduction

\documentclass[../../../include/open-logic-section]{subfiles}

\begin{document}

\olfileid[fr]{sol}{syn}{int}

\olsection{Introduction}

En logique du premier ordre, nous combinons les
symboles non logiques d'un langage donné, à savoir
ses !!{constant}s, ses !!{function}s et ses
!!{predicate}s, avec les symboles logiques pour
exprimer des propriétés des !!{structure}s du
premier ordre. La notion de satisfaction permet de
le faire : elle met en relation une
!!{structure}~$\Struct{M}$, une assignation de
!!{variable}s~$s$ et une !!{formula}~$!A$.
On a $\Sat{M}{!A}[s]$ si et seulement si ce
qu'exprime~$!A$ est vrai lorsque ses
!!{constant}s, !!{function}s et !!{predicate}s
reçoivent l'interprétation donnée par~$\Struct{M}$,
et ses !!{variable}s libres celle donnée par~$s$.
L'interprétation de l'!!{identity}~$\eq$ est
fixée par la définition de $\Sat{M}{!A}[s]$,
tout comme celle de~$\lforall$ et de~$\lexists$.
L'!!{identity} est toujours interprétée par la
relation d'identité sur le !!{domain}~$\Domain{M}$
de la structure, et les quantificateurs parcourent
toujours tout ce !!{domain}. Mais la quantification
ne porte, de manière décisive, que sur les éléments
du !!{domain}; seules des !!{variable}s d'objet
peuvent donc suivre un quantificateur.

En logique du second ordre, on étend le langage
et la définition de la satisfaction pour inclure
des variables de fonction et de prédicat, libres
ou liées, ainsi que leur quantification. Ces
variables sont aux !!{function}s et aux
!!{predicate}s ce que les variables d'objet sont
aux !!{constant}s. Elles jouent un rôle analogue
dans la formation des termes et des !!{formula}s
du second ordre, et l'on traite leur quantification
de façon analogue. Dans la \emph{sémantique
standard}, les quantificateurs du second ordre
parcourent tous les objets possibles du type voulu :
les fonctions à $n$ arguments de $\Domain{M}$
dans~$\Domain{M}$ pour les variables de fonction,
et les relations $n$-aires sur $\Domain{M}$ pour
les variables de prédicat.

Par exemple, la !!{formula}
\[
\lforall[\Obj{v_0}][(\Obj{P^1_0}(\Obj{v_0})
\lor\lnot\Obj{P^1_0}(\Obj{v_0}))]
\]
appartient aussi bien à la logique du premier
ordre qu'à celle du second ordre. Dans cette
dernière, nous pouvons également considérer
\[
\lforall[\Obj{V^1_0}][\lforall[\Obj{v_0}]
[(\Obj{V^1_0}(\Obj{v_0})\lor
\lnot\Obj{V^1_0}(\Obj{v_0}))]]
\]
et
\[
\lexists[\Obj{V^1_0}][\lforall[\Obj{v_0}]
[(\Obj{V^1_0}(\Obj{v_0})\lor
\lnot\Obj{V^1_0}(\Obj{v_0}))]].
\]
Ces deux dernières !!{formula}s n'ont aucune
!!{variable} libre : ce sont donc des
!!{sentence}s du second ordre. Ici,
$\Obj{V^1_0}$ est une variable de prédicat
unaire du second ordre. Ses interprétations
possibles sont les mêmes que celles d'un
!!{predicate} unaire tel que~$\Obj{P^1_0}$,
c'est-à-dire les sous-ensembles de~$\Domain{M}$.
Quantifier sur $\Obj{V^1_0}$ revient donc à
affirmer que
$\lforall[\Obj{v_0}]
[(\Obj{V^1_0}(\Obj{v_0})\lor
\lnot\Obj{V^1_0}(\Obj{v_0}))]$
vaut pour toutes les façons d'attribuer à
$\Obj{V^1_0}$ un sous-ensemble de~$\Domain{M}$,
ou pour au moins l'une d'elles. Puisque tout
ensemble contient ou ne contient pas un objet
donné, les deux !!{sentence}s sont vrais dans
toute !!{structure}.

\begin{editorial}
Dans la dernière occurrence de la !!{formula}
quantifiée, la source écrit une fois
$\Obj v_0$ et une fois $v_0$ pour la même
!!{variable} d'objet. Nous écrivons
$\Obj{v_0}$ dans les deux cas, comme dans les
!!{formula}s affichées juste avant.
\end{editorial}

\end{document}
